% English supplement, embedded-inclusion version; values in guillemets are filled by fill_embedded.py.
\documentclass[11pt,a4paper]{article}
\usepackage[T1]{fontenc}
\usepackage[utf8]{inputenc}
\usepackage[english]{babel}
\usepackage[margin=25mm]{geometry}
\usepackage{amsmath,amssymb,graphicx}
\usepackage{placeins}
\usepackage[numbers,sort&compress]{natbib}
\usepackage[hyphens]{url}
\usepackage[hidelinks]{hyperref}
\setlength{\emergencystretch}{2em}
\ifx\pdfmapfile\undefined\else
\pdfmapfile{+cm-super-t1.map}\pdfmapfile{+cm-super-ts1.map}\fi
\renewcommand{\thesection}{S\arabic{section}}
\renewcommand{\theequation}{S\arabic{equation}}
\renewcommand{\thefigure}{S\arabic{figure}}
\renewcommand{\thetable}{S\arabic{table}}

\title{Supplementary material to ``Stresses in aluminium around an
Al$_{13}$Fe$_4$ inclusion with a given magnetostrictive strain''}
\author{I.~Mikhailovskiy, D.~Pshonkin}
\date{}

\begin{document}
\maketitle

\section{Building and preparing the cell}

The matrix is a periodic aluminium crystal with lattice parameter
$a_0=4.05$~\AA{}, at which bulk aluminium is stress free with this potential.
The cell has 44 periods of 2.864~\AA{} along $x\parallel[1\bar10]$, 25 periods
of 4.960~\AA{} along $y\parallel[11\bar2]$ and 54 $(111)$ layers of 2.338~\AA{}
along $z\parallel[111]$; its size is $126.0\times124.0\times126.3$~\AA{}. The inclusion is a
ball of radius 25~\AA{} cut from an Al$_{13}$Fe$_4$ crystal with its $a$ axis
along $x$ and $b$ axis along $y$. It contains 4\,487 atoms, 3\,429 Al and
1\,058 Fe (ratio 3.24 against the stoichiometric 3.25). The aluminium atoms
inside the ball and all those closer than 2.2~\AA{} to an inclusion atom were
removed (4\,426 atoms). The shortest distances after construction are
2.20~\AA{} between matrix and inclusion and 2.37~\AA{} inside the inclusion.

With this potential a freely relaxed Al$_{13}$Fe$_4$ crystal changes its
lattice with respect to the measured one: $a$ by $+7.5\%$, $b$ by $-2.5\%$,
the cell height by $+3.0\%$, the volume by $+7.9\%$. The inclusion
atoms are therefore fixed at the measured lattice at every stage. In a trial run with a
free inclusion the cell stretched unevenly along its axes at 300~K, which
points to the inclusion swelling inside the aluminium; that run was stopped
and is not used.

Preparation of the reference cell: minimisation of the matrix energy to a
force two-norm of 1~eV/\AA{}; 10~ps at 300~K (Nos\'e--Hoover thermostat,
matrix atoms only, fixed cell size); cooling to 0.5~K over 10~ps (Langevin
thermostat); minimisation in chunks of 2000 iterations to a force two-norm of
$3\times10^{-4}$~eV/\AA{}. The mean normal stress of the matrix atoms more than
15~\AA{} from the inclusion surface is then $\sigma_0=«CTL_FAR_MEAN»$~MPa.
All cell dimensions are multiplied by $1-\sigma_0/3K=«CTL_BOX_SCALE»$, where
$K=77.4$~GPa is the bulk modulus of aluminium with the potential (only matrix
atoms are moved), and the minimisation is repeated. The mean stress far from
the inclusion is then $«CTL_FAR_AFTER»$~MPa.

\section{Convergence and linearity}

«S2_TEXT»

\section{Given strain and free magnetostriction}

The inclusion atoms are fixed, so the strain (1) of the main text is
the actual strain of the inclusion inside the matrix. The matrix field depends
only on the displacement of the inclusion surface and on the elasticity of the
matrix. Consider an elastic spherical particle with shear modulus $k\mu$,
where $\mu$ is the shear modulus of aluminium, and free magnetostriction
$\lambda$ of the same form as (1), that is, the strain the particle would have
in the field without the matrix. Its actual strain inside the matrix is
smaller than $\lambda$. Eshelby's solution for an inhomogeneous inclusion, one
whose elastic constants differ from those of the matrix
\cite{Eshelby1957,Mura1987}, gives for a deviatoric strain (one that does not
change the volume, like (1))
\begin{equation}
\varepsilon=A\lambda,\qquad A=\frac{\beta k}{1+\beta(k-1)},\qquad
\beta=\frac{2(4-5\nu)}{15(1-\nu)} ,
\end{equation}
where $\nu$ is Poisson's ratio (taken equal for inclusion and matrix for
simplicity); for $\nu=0.344$, $\beta=0.463$. For a given $\lambda$ all matrix
stresses equal the computed ones multiplied by $A\lambda/\varepsilon$. For an
inclusion as stiff as the matrix ($k=1$) $A=0.46$; for inclusions 2 and 2.5 times
stiffer, $A=0.63$ and $0.68$; for a rigid inclusion $A=1$. This
potential gives the elastic constants of Al$_{13}$Fe$_4$ at the equilibrium
lattice it predicts for this compound, which is overestimated (Section S1), so
they are not used to estimate $k$.

\section{Elastic solutions}

\paragraph{Single sphere.} A spherical inclusion of radius $R$ with a uniform
transformation strain $\boldsymbol\varepsilon^*$ (the strain it would take if
it were removed from the matrix) in an infinite isotropic matrix
has a uniform strain $\boldsymbol\varepsilon=\mathbf S:\boldsymbol\varepsilon^*$
inside, where $\mathbf S$ is the Eshelby tensor \cite{Eshelby1957,Mura1987}. To
make the interior strain equal the given (1), the transformation strain is set to
$\boldsymbol\varepsilon^*=\mathbf S^{-1}:\boldsymbol\varepsilon$; for a
deviatoric strain this is $\boldsymbol\varepsilon/\beta$. The strain outside,
$\varepsilon_{ij}(\mathbf x)=D_{ijkl}(\mathbf x)\,\varepsilon^*_{kl}$, follows
from the exterior Eshelby tensor \cite{JuSun1999}; with $\rho=R/r$ and
$\mathbf n=\mathbf x/r$
\begin{multline}
30(1-\nu)\rho^{-3}D_{ijkl}=(3\rho^2+10\nu-5)\delta_{ij}\delta_{kl}
+(3\rho^2-10\nu+5)(\delta_{ik}\delta_{jl}+\delta_{il}\delta_{jk})
+15(1-\rho^2)\delta_{ij}n_kn_l\\
+15(1-2\nu-\rho^2)\delta_{kl}n_in_j
+15(\nu-\rho^2)(\delta_{ik}n_jn_l+\delta_{jl}n_in_k+\delta_{il}n_jn_k+\delta_{jk}n_in_l)
+15(7\rho^2-5)n_in_jn_kn_l .
\end{multline}
The stress is $\sigma_{ij}=\lambda_L\varepsilon_{kk}\delta_{ij}+2\mu\varepsilon_{ij}$
with the isotropic constants of aluminium for the potential ($\mu=26.9$~GPa,
$\nu=0.344$, Voigt--Reuss--Hill average of $C_{11}$, $C_{12}$, $C_{44}$). The
program was checked on a centre of dilatation
($\boldsymbol\varepsilon^*=10^{-3}\mathbf I$): at $r=1.5R$ it gives
$\sigma_{rr}=-21.77$ and $\sigma_{\theta\theta}=10.89$~MPa, in agreement with Love's exact
solution.

\paragraph{The same periodic cell.} For a periodic array of inclusions with the
cell period the equilibrium equations are solved in Fourier space
\cite{MoulinecSuquet1998}. For a homogeneous medium of stiffness $\mathbf C$
(the cubic constants of aluminium for the potential, rotated to the cell axes)
and a transformation strain field $\chi(\mathbf x)\boldsymbol\varepsilon^*$, with
$\chi$ equal to one in the inclusion,
\begin{equation}
\hat\varepsilon_{ij}(\mathbf k)=\tfrac12\bigl[\xi_iN_{jm}(\boldsymbol\xi)
+\xi_jN_{im}(\boldsymbol\xi)\bigr]\tau_{ml}\xi_l\,\hat\chi(\mathbf k),
\qquad \mathbf N=(\boldsymbol\xi\cdot\mathbf C\cdot\boldsymbol\xi)^{-1},
\quad \boldsymbol\tau=\mathbf C:\boldsymbol\varepsilon^*,
\end{equation}
with $\boldsymbol\xi=\mathbf k/|\mathbf k|$; the $\mathbf k=0$ term is left out,
so the mean strain of the cell is zero, as in the fixed-size calculation. The
grid spacing is 1~\AA{} and the inclusion boundary is smoothed over one grid
cell. $\boldsymbol\varepsilon^*$ is adjusted until the mean strain inside the
inclusion equals the given (1): in the periodic cell it has to be 1.07--1.11
times larger, depending on the component, than for the single sphere. The same calculation gives the displacements
used as the starting guess for the matrix atoms in the $+\varepsilon$ cell. The
mean added stress in the matrix at fixed cell size is known exactly,
$\langle\Delta\boldsymbol\sigma\rangle=-\mathbf C:\boldsymbol\varepsilon\,f/(1-f)$;
for $f=3.3\%$ it is $«MEAN_PRED»$~MPa.

\section{The creep model}

A dislocation passes an obstacle by a thermal fluctuation at a rate
$\omega_0\exp(-\Delta G/k_{\mathrm B}T)$, where $\omega_0$ is the attempt frequency
and $\Delta G$ the height of the energy barrier. A shear stress $\tau$ does work $V^*\tau$
along the path over the barrier, so it lowers the barrier for forward jumps and
raises it for backward ones: $\Delta G_\pm=\Delta G_0\mp V^*\tau$
\cite{KocksArgonAshby1975,CaillardMartin2003}. If stresses of both signs occur
with equal weights in the volume around the inclusions, the mean rate factor is
$\tfrac12[\exp(x)+\exp(-x)]=\cosh x$, $x=V^*\tau/k_{\mathrm B}T$. For
dilute inclusions whose fields do not overlap, the relative increase
of the mean rate in a material with an inclusion volume fraction $f$ is the
integral of $\cosh x-1$ over the matrix around one inclusion times the number
of inclusions per unit volume, $f/V_p$; this is Eq.~(2) of the main text. For a
field $\tau(r)=\tau_m(R/r)^3$, where $\tau_m$ is the shear stress at the
inclusion surface, the substitution $u=x_0(R/r)^3$ reduces the
integral to
\begin{equation}
H=f\,x_0\int_0^{x_0}\frac{\cosh u-1}{u^2}\,\mathrm du,\qquad
x_0=\frac{V^*\tau_m}{k_{\mathrm B}T},
\end{equation}
which gives the solid curve of Fig.~5. The dashed curve uses the computed field
of the single sphere: the integral (2) is taken on a spherical grid (120 radial
nodes evenly spaced in $(R/r)^3$, which cover the whole matrix, and
$48\times48$ angular nodes) with the shear stresses on the
twelve $\{111\}\langle110\rangle$ slip systems for 64 random grain orientations.
The amplitude $\tau_m$ for which Eq.~(S4) gives the same $H$ is then found.
A 25\% increase is reached at $V^*=«VSTAR_GRAINS»\,b^3$ (average over grain
orientations), «VSTAR_CRYSTAL»$\,b^3$ (one orientation, that of the cell) and
«VSTAR_MAX»$\,b^3$ (for the same orientation, the largest of the shear stresses
on the twelve slip systems at every point). Since the field is proportional to
$\varepsilon$, $H$ depends only on the product $V^*\varepsilon$, and at
$V^*=70\,b^3$ a 25\% increase needs $\varepsilon\cdot«VSTAR_GRAINS»/70\approx
«EPS_AT_70»\%$.

\section{Lattice resistance to dislocation glide}

To judge whether the motion of a dislocation past the inclusion can be
followed atomistically, the stress at which a straight edge dislocation
glides with this potential was measured in a cell without an inclusion. The
cell is an aluminium slab of $126.0\times19.8\times125.5$~\AA{} (18\,768
atoms), periodic along $x$ and $y$ with free surfaces along $z$; the three
top and three bottom $(111)$ layers are rigid. An edge dislocation with
Burgers vector $b=2.864$~\AA{} along $x$ and line along $y$ is put into the
slip plane at a height of 91.2~\AA{}: the upper half of the slab has one
more column of atoms along $x$ than the lower half. The top block is moved
along $x$ in steps of 0.00455~\AA, and after each step the energy is
minimised to a force norm of $10^{-3}$~eV/\AA{} at 0~K. The applied shear
stress is the force of the matrix on the top and bottom blocks divided by
the area $L_xL_y$ (the mean of the two blocks, which differ by at most
1~MPa).

The stress rises linearly while the dislocation stays put and drops when it
jumps into the next valley of the lattice relief, that is, the next position
of lowest dislocation energy. The valleys are
1.41--1.42~\AA{} apart, that is $b/2$; the line positions found from the
block displacement agree with those found by DXA \cite{Stukowski2012DXA}
to within 0.25~\AA, and DXA finds two partial dislocations in every frame:
in aluminium an edge dislocation splits into two partials separated by a
ribbon where the stacking of the $(111)$ planes is faulted.
Before each of the three jumps the stress reaches 9.13, 9.00 and 9.01~MPa,
so with this potential a straight edge dislocation at 0~K starts to glide at
$\tau_P\approx9$~MPa (the Peierls stress).

«S6_TEXT»

\section{Data and programs}
\sloppy

Programs and inputs (paths from the repository root):
\begin{itemize}
\item the cell: \path{scripts/build_embedded_inclusion.py}; the strained
cells: \path{scripts/make_strained_inclusion.py}; the LAMMPS inputs:
\path{lammps/embedded_inclusion/in.relax_control} and
\path{lammps/embedded_inclusion/in.static_pair};
\item stresses: \path{analysis/python/embedded_inclusion_field.py};
convergence and linearity checks: \path{analysis/python/embedded_gates.py};
defect search: \path{analysis/python/embedded_dxa.py}; strain, stress per
unit strain and density at the interface:
\path{analysis/python/embedded_interface.py};
elastic solutions: \path{analysis/python/embedded_inclusion_continuum.py};
\item the creep model: \path{analysis/python/embedded_creep.py}; the yield
threshold at the particle: \path{analysis/python/embedded_yield.py}; elastic
constants of the potential: \path{analysis/python/meam_elastic_constants.py};
figures: \path{analysis/python/embedded_figures.py};
\item the dislocation slab: \path{scripts/build_embedded_dislocation.py};
relaxation: \path{lammps/embedded_inclusion/in.disl_relax}; dislocation
positions: \path{analysis/python/embedded_disl_states.py}; shear:
\path{lammps/embedded_inclusion/in.disl_sweep} and
\path{analysis/python/embedded_sweep.py};
\item the runs of the $\pm\varepsilon$ cells: \path{scripts/queue_embedded_static.sh};
the runs after them (the dislocation above the inclusion, the cells with
strains $4\varepsilon$ and $10\varepsilon$, shear): \path{scripts/queue_embedded_plasticity.sh};
records of the dislocation runs and of the $4\varepsilon$ and $10\varepsilon$
cells: \path{analysis/python/embedded_records.py}.
\end{itemize}
Result records: \path{docs/reports/embedded_field.json},
\path{docs/reports/embedded_gates.json}, \path{docs/reports/embedded_continuum.json},
\path{docs/reports/embedded_creep.json}, \path{docs/reports/embedded_yield.json},
\path{docs/reports/embedded_peierls.json}, \path{docs/reports/embedded_dxa.json},
\path{docs/reports/embedded_interface.json}, \path{docs/reports/embedded_margin.json},
\path{docs/reports/embedded_disl_states.json},
\path{docs/reports/meam_elastic_constants.json}. The configurations and run
logs, compressed and with checksums, are in \path{data/publication/embedded}.

\bibliographystyle{unsrtnat}
\bibliography{references}
\end{document}
